% =============================================================================
% §6 Interpretation (R5) — SDT channel-assisted compression; cohort-diffusion ↔
% cohort-dominance; Esteve anchor; ARG/CHL falsifiable prediction.
% Non-load-bearing: the paper stands on R1–R4.
% =============================================================================
\section{Interpretation: Channel-Assisted Compression of the Second
Demographic Transition}
\label{sec:interpretation}

The results so far stand on their own: a measured collapse
(Section~\ref{sec:phenomenon}), an identified structure
(Section~\ref{sec:discrimination}), and a sufficient parsimonious mechanism
(Section~\ref{sec:model}). This section offers the frame that makes them
cohere; the paper's claims do not depend on it.

The Second Demographic Transition describes ideational change ---
individualization, secularization, the decoupling of partnership from marriage
and of marriage from childbearing --- diffusing across cohorts
\citep{lesthaeghe1986twee, vandekaa1987, lesthaeghe2010}. Its standard account
explains direction well and tempo poorly: cohort-borne ideational change is
ordinarily a slow mechanism. What Latin America adds is a pre-existing
\emph{channel}. The region has maintained, since the colonial period, a dual
nuptiality system in which consensual union is a socially recognized,
institutionally familiar alternative to marriage --- massively expanded in the
recent ``cohabitation boom'' \citep{esteve2012}. Where that channel is already
wide, an arriving cohort that has absorbed SDT preferences does not need to
pioneer a new living arrangement, defy strong sanctions, or wait for
institutions to adapt; it flows into an existing one. The channel converts
ideational change into behavioral change at the speed of cohort succession
rather than the slower speed of institutional innovation --- \emph{compression},
assisted by the channel.

This reading knits the empirical pieces together. It explains why the change
localizes on cohort lines with stabilizing within-cohort dynamics
(Section~\ref{sec:discrimination}): the action is between cohorts, not within
them. It explains why a three-parameter cohort gradient suffices
(Section~\ref{sec:model}): the mechanism genuinely is a compressed
across-cohort shift in marriage propensity. And it explains the convergence of
our two independent lenses --- the econometric cohort \emph{dominance} and the
demographic account of cohort \emph{diffusion} are two faces of the same
object, and the model's recovered inflection at birth cohorts 1989--1992 dates
where the compression concentrated.

The frame also does what a frame should: it generates a falsifiable
out-of-sample prediction. If the consensual-union channel is the speed
mechanism, then in marriage-dominant regimes where the channel is narrow ---
Argentina and Chile are the natural regional cases, with the channel-relevant
composition series still under construction --- fertility decline of comparable
depth should be \emph{slower or differently structured}: less
marriage-to-cohabitation substitution at flat union totals, more decline within
existing union types, and weaker cohort localization. Both countries have
experienced deep declines (Section~\ref{sec:phenomenon}); the composition
anatomy of those declines is the test. We state the prediction now, before the
data are assembled, in the same pre-registration spirit as the rest of the
paper. Confirmation would establish the channel as the compression mechanism;
disconfirmation would bound the frame without touching the identified
cohort-replacement structure of Sections~\ref{sec:discrimination}
and~\ref{sec:model}.
